\chapter{Mô hình hệ cơ sở tri thức (KBS)}

\section{Kiến trúc logic}
Engine lớp \texttt{KnowledgeRuleEngine} (\texttt{kbs/knowledge\_rules.py}) nạp cấu hình \texttt{rules\_config.json}, phiên bản 3.0, phân nhánh \texttt{khtn\_rules} và \texttt{khxh\_rules}. Mỗi ngành là một đối tượng chứa danh sách luật; mỗi luật có các trường: \texttt{name}, \texttt{thresholds}, \texttt{operator}, \texttt{score}, \texttt{specificity}, \texttt{reason} \cite{docs_rules}.

Luồng xử lý cho một vector điểm và một chỉ số ngành:
\begin{enumerate}[leftmargin=*]
  \item Duyệt tất cả luật của ngành, kiểm tra điều kiện.
  \item Thu thập tập luật khớp; nếu rỗng, gán điểm mặc định thấp (theo implementation).
  \item Áp dụng \texttt{resolve\_conflicts}: ưu tiên \texttt{specificity} cao nhất, sau đó \texttt{score}.
  \item (Tuỳ chọn) áp dụng chuỗi suy luận \texttt{forward\_chain} nếu cấu hình \texttt{chaining\_rules} khớp ngữ cảnh.
\end{enumerate}

\section{Toán tử điều kiện}
Hai toán tử được dùng:
\begin{itemize}[leftmargin=*]
  \item \texttt{AND}: tất cả các ràng buộc dạng ``điểm môn $\geq$ ngưỡng'' phải đồng thời đúng.
  \item \texttt{OR\_LESS\_THAN}: dùng cho luật \textit{không phù hợp} --- khi \textbf{bất kỳ} điều kiện dạng điểm $<$ ngưỡng xảy ra thì luật kích hoạt (điểm KBS thấp, mang ý nghĩa cảnh báo).
\end{itemize}

\begin{algorithm}[!htbp]
\caption{Giải quyết xung đột luật (chiến lược \texttt{specificity\_then\_score})}
\label{alg:conflict}
\begin{algorithmic}[1]
\Require Điểm thí sinh $s$; tập luật khớp $\mathcal{R}$ (mỗi luật có \texttt{specificity}, \texttt{score})
\State Sắp xếp $\mathcal{R}$ theo cặp $(\texttt{specificity},\texttt{score})$ giảm dần thứ tự từ điển
\State \textbf{return} luật đứng đầu danh sách và điểm tương ứng
\end{algorithmic}
\end{algorithm}
\FloatBarrier

\section{Bảng tóm tắt luật theo ngành (KHTN)}
Mỗi ngành có bốn mức: Rất phù hợp, Phù hợp, Trung bình, Không phù hợp. Bảng \ref{tab:khtnrules} tóm tắt ngưỡng chính (chi tiết đầy đủ nằm trong JSON).

\begin{table}[!htbp]
\centering
\caption{Tóm tắt luật KHTN (trích ngưỡng chính)}
\label{tab:khtnrules}
\small
\begin{tabular}{p{2.2cm}p{2.5cm}p{2.2cm}p{2.2cm}p{2.2cm}c}
\toprule
Ngành & Rất phù hợp & Phù hợp & Trung bình & Không phù hợp & Điểm max \\
\midrule
IT & To$\geq8$, Lý$\geq7{,}5$, Anh$\geq6$ & To$\geq7$, Lý$\geq6{,}5$, Hóa$\geq5$, Anh$\geq5$ & To$\geq7$, Lý$\geq6$ & To$<6$ hoặc Lý$<5{,}5$ & 95 \\
Y khoa & Sinh$\geq8{,}5$, Hóa$\geq8$, Lý$\geq7$ & Sinh$\geq8$, Hóa$\geq7{,}5$, Lý$\geq6$, Văn$\geq6$ & Sinh$\geq7{,}5$, Hóa$\geq7$ & Sinh$<6{,}5$ hoặc Hóa$<6$ & 95 \\
Kinh tế & Anh$\geq8$, To$\geq7{,}5$, Văn$\geq7$ & Anh$\geq7$, To$\geq6{,}5$, Văn$\geq6{,}5$ & Anh$\geq6{,}5$, To$\geq6$ & Anh$<6$ hoặc To$<5{,}5$ & 90 \\
Kỹ thuật & To$\geq8$, Lý$\geq8$, Hóa$\geq7$ & To$\geq7{,}5$, Lý$\geq7$, Hóa$\geq6{,}5$ & To$\geq7{,}5$, Lý$\geq6{,}5$ & To$<6{,}5$ hoặc Lý$<6$ & 92 \\
Nông--Lâm--Ngư & Sinh$\geq8$, Hóa$\geq7{,}5$, To$\geq6$ & Sinh$\geq7{,}5$, Hóa$\geq7$, To$\geq5{,}5$ & Sinh$\geq7$, Hóa$\geq6{,}5$ & Sinh$<6{,}5$ hoặc Hóa$<5{,}5$ & 88 \\
\bottomrule
\end{tabular}
\end{table}
\FloatBarrier

\section{Bảng tóm tắt luật theo ngành (KHXH)}
Bảng \ref{tab:khxhrules} trích ngưỡng từ \texttt{rules\_config.json} (nhánh \texttt{khxh\_rules}). Sử $=$ \texttt{lich\_su}, Địa $=$ \texttt{dia\_ly}. Cột \textit{Không phù hợp} dùng \texttt{OR\_LESS\_THAN}.

\begin{table}[!htbp]
\centering
\caption{Tóm tắt luật KHXH theo \texttt{rules\_config.json} (v3.0)}
\label{tab:khxhrules}
\footnotesize
\setlength{\tabcolsep}{3pt}
\begin{tabular}{p{1.9cm}p{2.35cm}p{2.35cm}p{2.1cm}p{2.05cm}c}
\toprule
Ngành & Rất phù hợp & Phù hợp & Trung bình & Không phù hợp & Điểm \\
 & \textit{(AND)} & \textit{(AND)} & \textit{(AND)} & \textit{(OR $<$)} & \textit{max} \\
\midrule
Kinh tế &
Anh$\geq8$, To$\geq7{,}5$, Văn$\geq7$ &
Anh$\geq7$, To$\geq6{,}5$, Văn$\geq6{,}5$ &
Anh$\geq6{,}5$, To$\geq6$ &
Anh$<6$ \textbf{hoặc} To$<5{,}5$ & 90 \\
\addlinespace[2pt]
Sư phạm &
Văn$\geq8$, Anh$\geq7{,}5$, Sử$\geq7$ &
Văn$\geq7$, Anh$\geq7$, Sử$\geq6{,}5$ &
Văn$\geq6{,}5$, Anh$\geq6{,}5$ &
Văn$<6$ \textbf{hoặc} Anh$<5{,}5$ & 90 \\
\addlinespace[2pt]
Luật &
GDCD$\geq8$, Sử$\geq8$, Văn$\geq7$ &
GDCD$\geq7$, Sử$\geq7$, Văn$\geq6{,}5$, Anh$\geq6$ &
GDCD$\geq6{,}5$, Sử$\geq6{,}5$ &
GDCD$<6$ \textbf{hoặc} Sử$<5{,}5$ & 90 \\
\addlinespace[2pt]
Du lịch &
Anh$\geq8$, Địa$\geq7{,}5$, Văn$\geq8$ &
Anh$\geq7{,}5$, Địa$\geq7$, Văn$\geq6{,}5$ &
Anh$\geq6{,}5$, Địa$\geq6{,}5$, Văn$\geq6$ &
Anh$<6{,}5$ \textbf{hoặc} Địa$<5{,}5$ & 88 \\
\bottomrule
\end{tabular}
\end{table}

\noindent\textbf{Điểm số (\texttt{score}) bốn mức} (0--100, sau bước ưu tiên \texttt{specificity}): Kinh tế 90\,/\,75\,/\,55\,/\,15; Sư phạm 90\,/\,75\,/\,60\,/\,15; Luật 90\,/\,75\,/\,60\,/\,15; Du lịch 88\,/\,78\,/\,68\,/\,15. Ví dụ \texttt{Luat\_Fit} có \texttt{specificity}=4 (bốn ngưỡng AND).
\FloatBarrier

\section{Suy luận chuỗi (\textit{forward chaining})}
Ngoài luật tĩnh theo ngành, \texttt{rules\_config.json} có thể khai báo \texttt{chaining\_rules} theo khối. Sau khi chọn luật nền, \texttt{forward\_chain} duyệt các bước bổ sung điều chỉnh điểm hoặc gắn thẻ giải thích khi điều kiện phụ thỏa. Cơ chế này mở đường mở rộng sang luật đa tầng (ví dụ: nếu IT\_Fit và điểm Tin học nội bộ giả lập $>8$ thì tăng điểm) mà không phá vỡ kiến trúc JSON hiện tại.

\section{Giải thích và minh họa xung đột}
Giả sử học sinh KHTN có Toán 8{,}5, Lý 8{,}0, Anh 5{,}5: luật IT\_Very\_Fit yêu cầu Anh $\geq 6$ nên không còn hiệu lực đầy đủ; luật IT\_Fit có thể khớp với bốn điều kiện và \texttt{specificity}=4, có thể được chọn thay vì luật ba điều kiện tùy tập khớp --- minh họa tinh thần ưu tiên luật chi tiết hơn \cite{docs_rules}.

\section{Điểm mạnh và hạn chế của KBS trong dự án}
\textbf{Điểm mạnh:} cập nhật luật không cần biên dịch lại toàn bộ mã (chỉ cần kiểm tra JSON hợp lệ); mỗi kết quả gắn chuỗi \texttt{reason} phục vụ báo cáo giải thích.

\textbf{Hạn chế:} ngưỡng cứng; chưa có lớp học từ dữ liệu cho luật; chưa có quy trình phê duyệt chuyên gia giáo dục bắt buộc.

\section{Kết chương}
KBS đóng vai trò ràng buộc miền và giải thích; số hóa điểm từ 0--100 đồng bộ với nhánh học máy trước khi vào fusion.
